\section{The Appointed Texts}
\label{sec:appointed-texts}

The Latin below is that of the 1962 typical edition, \work{Missale Romanum}
(Typis Polyglottis Vaticanis), printed pp.~420--421, nos.~1749--1758, read on
the page images of the CMAA facsimile. The Missal prints the three orations,
the Epistle and the Gospel a second time in the season after the Epiphany
(pp.~47--48), with the same words; the Latin here follows the resumed
printing, in which the Collect has its conclusion written out and the Secret
and Postcommunion have the cue \latin{Per Dóminum.} Every word of all ten
texts stands in the Ratisbon printing of the \work{Missale Romanum} (Pustet,
1862): the orations and lessons under the Fifth Sunday after the Epiphany,
pp.~42--44, and the chants under the Twenty-third Sunday after Pentecost,
pp.~357--358, where a rubric directs that the resumed Sundays take those
chants. The two printings differ in spelling (\latin{Jesus}, \latin{coeléstis},
\latin{charitátem}), in pointing and in the length of conclusion cues, and in
no word that is said.

The Lasance \work{New Roman Missal} in the Benziger revision of 1945 prints
this whole Mass, Latin and English, as its \work{Supplementary Mass No.~2},
``Mass of the Fifth Sunday after Epiphany,'' pp.~744--747, under the rubric
that it is said on the twenty-fourth Sunday after Pentecost when there are
twenty-six. Its English of the five chants renders the Missal's own forms of
the verses where they differ from the Clementine Vulgate, and that English
is given for each chant. The Douay--Rheims Bible in Bishop
Challoner's revision, which translates the Clementine Vulgate, is given for
the two lessons and beside each chant at the canonical verses (for the
Offertory, which repeats the Alleluia's verses, under the Alleluia), and each place
where the Missal's form departs from the Vulgate is named beneath the text.
The English of the Collect and Postcommunion is the Lasance's under the Fifth
Sunday after the Epiphany, pp.~196 and 198, and of the Secret its
supplementary Mass, p.~746. None of this English is an approved liturgical
translation of the 1962 Missal. Psalms are cited in the Missal's numbering
throughout; the Hebrew numbering is added in parentheses at each appointed
psalm text below and in the map of the texts.

\subsection{Introit}
\label{proper-introit}

\properlocus{\latin{Antiphona ad Introitum}, \latin{Ierem.\ 29, 11, 12 et 14};
verse \latin{Ps.\ 84, 2} (Heb.\ 85). Missal p.~420, no.~1749.}

\begin{latinproper}
Dicit Dóminus: Ego cógito cogitatiónes pacis, et non afflictiónis:
invocábitis me, et ego exáudiam vos: et redúcam captivitátem vestram de
cunctis locis. \textnormal{Ps.\ 84, 2} Benedixísti, Dómine, terram tuam:
avertísti captivitátem Iacob. ℣.~Glória Patri.
\end{latinproper}

\begin{englishwitness}{Lasance, The New Roman Missal (1945), Supplementary Mass No.~2, Introit}
The Lord saith, I think thoughts of peace, and not of affliction: you shall
call upon Me, and I will hear you; and I will bring back your captivity from
all places. Ps.~84, 2. Lord, Thou hast blessed Thy land: Thou hast turned away
the captivity of Jacob. ℣.~Glory.
\end{englishwitness}

\begin{englishwitness}{Douay--Rheims (Challoner), Jer 29:11--12, 14, the antiphon's verses}
For I know the thoughts that I think towards you, saith the Lord, thoughts of
peace, and not of affliction, to give you an end and patience. And you shall
call upon me, and you shall go: and you shall pray to me, and I will hear you.
\dots\ And I will be found by you, saith the Lord: and I will bring back your
captivity, and I will gather you out of all nations, and from all the places
to which I have driven you out \dots
\end{englishwitness}

\noindent The antiphon is not a continuous verse of the Vulgate. It opens
\latin{Dicit Dóminus}, where the Vulgate's verse~11 has God say \latin{Ego
enim scio cogitationes quas ego cogito super vos, ait Dominus}. It has nothing
of \latin{ut dem vobis finem et patientiam} (``to give you an end and
patience''), of \latin{et ibitis: et orabitis me} (``and you shall go: and you
shall pray to me''), of verse~13, of the opening of verse~14, \latin{Et
inveniar a vobis, ait Dominus} (``And I will be found by you, saith the
Lord''), or of \latin{et congregabo vos de universis gentibus} (``and I will
gather you out of all nations''), and it joins
\latin{redúcam captivitátem vestram} directly to \latin{de cunctis locis}. The
psalm verse agrees with the Vulgate. After the \latin{Gloria Patri} the
antiphon is repeated.

\subsection{Collect}
\label{proper-collect}

\properlocus{\latin{Oratio}. Missal p.~420, no.~1750.}

\begin{latinproper}
Famíliam tuam, quǽsumus, Dómine, contínua pietáte custódi: ut, quæ in sola
spe grátiæ cæléstis innítitur, tua semper protectióne muniátur. Per Dóminum
nostrum Iesum Christum, Fílium tuum: Qui tecum vivit et regnat in unitáte.
\end{latinproper}

\begin{englishwitness}{Lasance, The New Roman Missal (1945), Fifth Sunday after Epiphany, Collect}
Keep Thy family, we beseech Thee, O Lord, with Thy continual mercy that,
leaning only upon the hope of Thy heavenly grace, it may ever be defended by
Thy protection. Through our Lord.
\end{englishwitness}

\noindent The Lasance renders \latin{pietas} as ``mercy.'' The relative
clause \latin{quæ in sola spe grátiæ cæléstis innítitur}, ``which leans on the
hope of heavenly grace alone,'' stands between the petition and its purpose.

\subsection{Epistle}
\label{proper-epistle}

\properlocus{\latin{Léctio Epístolæ beáti Pauli Apóstoli ad Colossénses},
Col 3:12--17. Missal p.~420, no.~1751.}

\begin{latinproper}
Fratres: Indúite vos sicut elécti Dei, sancti et dilécti, víscera
misericórdiæ, benignitátem, humilitátem, modéstiam, patiéntiam: supportántes
ínvicem, et donántes vobismetípsis, si quis advérsus áliquem habet querélam:
sicut et Dóminus donávit vobis, ita et vos. Super ómnia autem hæc, caritátem
habéte, quod est vínculum perfectiónis: et pax Christi exsúltet in córdibus
vestris, in qua et vocáti estis in uno córpore: et grati estóte. Verbum Christi
hábitet in vobis abundánter, in omni sapiéntia, docéntes, et commonéntes
vosmetípsos, psalmis, hymnis, et cánticis spirituálibus, in grátia cantántes
in córdibus vestris Deo. Omne quodcúmque fácitis in verbo aut in ópere, ómnia
in nómine Dómini Iesu Christi, grátias agéntes Deo et Patri per Iesum Christum
Dóminum nostrum.
\end{latinproper}

\begin{englishwitness}{Douay--Rheims (Challoner), Col 3:12--17}
Put ye on therefore, as the elect of God, holy and beloved, the bowels of
mercy, benignity, humility, modesty, patience: bearing with one another and
forgiving one another, if any have a complaint against another. Even as the
Lord hath forgiven you, so do you also. But above all these things have
charity, which is the bond of perfection. And let the peace of Christ rejoice
in your hearts, wherein also you are called in one body: and be ye thankful.
Let the word of Christ dwell in you abundantly: in all wisdom, teaching and
admonishing one another in psalms, hymns and spiritual canticles, singing in
grace in your hearts to God. All whatsoever you do in word or in work, do all
in the name of the Lord Jesus Christ, giving thanks to God and the Father by
him.
\end{englishwitness}

\noindent The Missal opens \latin{Fratres: Indúite vos sicut}, where the
Vulgate reads \latin{Induite vos ergo, sicut}, ``Put ye on therefore''; and
it ends \latin{per Iesum Christum Dóminum nostrum}, ``through Jesus Christ our
Lord,'' where the Vulgate reads \latin{per ipsum}, ``by him.'' Every other
word of the six verses agrees.

\subsection{Gradual}
\label{proper-gradual}

\properlocus{\latin{Graduale}, Ps 43:8--9 (44:8--9). Missal p.~420,
no.~1752.}

\begin{latinproper}
Liberásti nos, Dómine, ex affligéntibus nos: et eos, qui nos odérunt,
confudísti. ℣.~In Deo laudábimur tota die, et in nómine tuo confitébimur in
sǽcula.
\end{latinproper}

\begin{englishwitness}{Lasance, The New Roman Missal (1945), Supplementary Mass No.~2, Gradual}
Thou hast delivered us, O Lord, from them that afflict us: and hast put them
to shame that hate us. ℣.~In God we will glory all the day: and in Thy name we
will give praise for ever.
\end{englishwitness}

\begin{englishwitness}{Douay--Rheims (Challoner), Ps 43:8--9}
But thou hast saved us from them that afflict us: and hast put them to shame
that hate us. In God shall we glory all the day long: and in thy name we will
give praise for ever.
\end{englishwitness}

\noindent The Vulgate's verse~8 reads \latin{salvasti enim nos de affligentibus
nos, et odientes nos confudisti}; the Missal sings \latin{Liberásti nos,
Dómine, ex affligéntibus nos: et eos, qui nos odérunt, confudísti}. Verse~9
ends \latin{in saeculum} in the Vulgate and \latin{in sǽcula} in the Missal.
The Latin psalm text that St Ambrose expounds reads \latin{Liberasti enim nos
ex affligentibus nos}, and so does the psalm text printed with the
commentary of Cassiodorus.

\subsection{Alleluia}
\label{proper-alleluia}

\properlocus{Ps 129:1--2 (130:1--2). Missal p.~420, no.~1753.}

\begin{latinproper}
Allelúia, allelúia. ℣.~De profúndis clamávi ad te, Dómine: Dómine, exáudi
oratiónem meam. Allelúia.
\end{latinproper}

\begin{englishwitness}{Lasance, The New Roman Missal (1945), Supplementary Mass No.~2, Alleluia}
Alleluia, alleluia. ℣.~From the depths I have cried to Thee, O Lord: Lord,
hear my prayer. Alleluia.
\end{englishwitness}

\begin{englishwitness}{Douay--Rheims (Challoner), Ps 129:1--2}
Out of the depths I have cried to thee, O Lord: Lord, hear my voice. \dots
\end{englishwitness}

\noindent The Vulgate reads \latin{Domine, exaudi vocem meam}, ``Lord, hear my
voice''; the Missal sings \latin{oratiónem meam}, ``my prayer.'' The verse
stops before the second half of verse~2, ``Let thy ears be attentive to the
voice of my supplication.'' There is no Tract and no Sequence.

\subsection{Gospel}
\label{proper-gospel}

\properlocus{\latin{Sequéntia sancti Evangélii secúndum Matthǽum}, Mt
13:24--30. Missal p.~421, no.~1754.}

\begin{latinproper}
In illo témpore: Dixit Iesus turbis parábolam hanc: Símile factum est regnum
cælórum hómini, qui seminávit bonum semen in agro suo. Cum autem dormírent
hómines, venit inimícus eius, et superseminávit zizánia in médio trítici, et
ábiit. Cum autem crevísset herba, et fructum fecísset, tunc apparuérunt et
zizánia. Accedéntes autem servi patrisfamílias, dixérunt ei: Dómine, nonne
bonum semen seminásti in agro tuo? Unde ergo habet zizánia? Et ait illis:
Inimícus homo hoc fecit. Servi autem dixérunt ei: Vis, imus, et collígimus
ea? Et ait: Non: ne forte colligéntes zizánia, eradicétis simul cum eis et
tríticum. Sínite útraque créscere usque ad messem, et in témpore messis dicam
messóribus: Collígite primum zizánia, et alligáte ea in fascículos ad
comburéndum, tríticum autem congregáte in hórreum meum.
\end{latinproper}

\begin{englishwitness}{Douay--Rheims (Challoner), Mt 13:24--30}
Another parable he proposed to them, saying: The kingdom of heaven is likened
to a man that sowed good seed in his field. But while men were asleep, his
enemy came and oversowed cockle among the wheat and went his way. And when the
blade was sprung up, and had brought forth fruit, then appeared also the
cockle. And the servants of the good man of the house coming said to him. Sir,
didst thou not sow good seed in thy field? Whence then hath it cockle? And he
said to them: An enemy hath done this. And the servants said to him: Wilt
thou that we go and gather it up? And he said: No, lest perhaps gathering up
the cockle, you root up the wheat also together with it. Suffer both to grow
until the harvest, and in the time of the harvest I will say to the reapers:
Gather up first the cockle, and bind it into bundles to burn, but the wheat
gather ye into my barn.
\end{englishwitness}

\noindent The Missal's opening, \latin{In illo témpore: Dixit Iesus turbis
parábolam hanc}, ``At that time Jesus spoke this parable to the crowds,''
stands for the Vulgate's \latin{Aliam parabolam proposuit illis, dicens}; every
other word of the seven verses agrees. The Credo follows the Gospel.

\subsection{Offertory}
\label{proper-offertory}

\properlocus{\latin{Ant.\ ad Offertorium}, Ps 129:1--2 (130:1--2). Missal
p.~421, no.~1755.}

\begin{latinproper}
De profúndis clamávi ad te, Dómine: Dómine, exáudi oratiónem meam: de
profúndis clamávi ad te, Dómine.
\end{latinproper}

\begin{englishwitness}{Lasance, The New Roman Missal (1945), Supplementary Mass No.~2, Offertory}
From the depths I have cried out to Thee, O Lord; Lord, hear my prayer: from
the depths I have cried out to Thee, O Lord.
\end{englishwitness}

\noindent The Offertory sings the same verses as the Alleluia, with the same
\latin{oratiónem meam}, and repeats the first line.

\subsection{Secret}
\label{proper-secret}

\properlocus{\latin{Secreta}. Missal p.~421, no.~1756.}

\begin{latinproper}
Hóstias tibi, Dómine, placatiónis offérimus: ut et delícta nostra miserátus
absólvas, et nutántia corda tu dírigas. Per Dóminum.
\end{latinproper}

\begin{englishwitness}{Lasance, The New Roman Missal (1945), Supplementary Mass No.~2, Secret}
We offer Thee, O Lord, the sacrifice of reconciliation, that Thou mayest
mercifully forgive our sins and direct our wavering hearts. Through our Lord.
\end{englishwitness}

\noindent \latin{Hostiae placationis} are offerings that appease or
propitiate; the Lasance renders the phrase ``the sacrifice of
reconciliation.'' After the Secret the Missal directs the Preface of the Most
Holy Trinity.

\subsection{Communion}
\label{proper-communion}

\properlocus{\latin{Ant.\ ad Communionem}, Mk 11:24. Missal p.~421,
no.~1757.}

\begin{latinproper}
Amen dico vobis, quidquid orántes pétitis, crédite quia accipiétis, et fiet
vobis.
\end{latinproper}

\begin{englishwitness}{Lasance, The New Roman Missal (1945), Supplementary Mass No.~2, Communion}
Amen, I say to you, whatsoever you ask when you pray, believe that you shall
receive, and it shall be done unto you.
\end{englishwitness}

\begin{englishwitness}{Douay--Rheims (Challoner), Mk 11:23--24}
Amen I say to you that whosoever shall say to this mountain, Be thou removed
and be cast into the sea, and shall not stagger in his heart, but believe that
whatsoever he saith shall be done; it shall be done unto him. Therefore I say
unto you, all things, whatsoever you ask when ye pray, believe that you shall
receive: and they shall come unto you.
\end{englishwitness}

\noindent The Vulgate's verse~24 reads \latin{Propterea dico vobis, omnia
quaecumque orantes petitis, credite quia accipietis, et evenient vobis}. The
antiphon begins with \latin{Amen dico vobis}, the words that open verse~23,
and ends \latin{et fiet vobis}, where verse~23 ends \latin{fiet ei}; between
them it gives verse~24's petition with \latin{quidquid} in place of the
Vulgate's \latin{omnia quaecumque}: \latin{quidquid orántes pétitis, crédite
quia accipiétis}.

\subsection{Postcommunion}
\label{proper-postcommunion}

\properlocus{\latin{Postcommunio}. Missal p.~421, no.~1758.}

\begin{latinproper}
Quǽsumus, omnípotens Deus: ut illíus salutáris capiámus efféctum, cuius per
hæc mystéria pignus accépimus. Per Dóminum.
\end{latinproper}

\begin{englishwitness}{Lasance, The New Roman Missal (1945), Fifth Sunday after Epiphany, Postcommunion}
We pray Thee, O almighty God, that we may receive the effect of that salvation
of which we have received the pledge in these mysteries. Through our Lord.
\end{englishwitness}

\noindent The Lasance's supplementary Mass, pp.~746--747, prints the same English
without ``in these mysteries''; the English of p.~198, given here, renders the
Latin \latin{per hæc mystéria}.
